\documentclass[12pt]{article}

% Idioma e codificação
\usepackage[brazil]{babel}
\usepackage[utf8]{inputenc}
\usepackage[T1]{fontenc}

% Matemática e listas
\usepackage{amsmath, amssymb, amsfonts}
\usepackage{enumitem}
\usepackage{tabularx}
\usepackage{booktabs}

% Layout e cores
\usepackage{geometry}
\usepackage{graphicx}
\usepackage{xcolor}

% Cores institucionais
\definecolor{maincolor}{RGB}{0,100,50}
\definecolor{forestgreen}{RGB}{0,100,50}

% Cabeçalho e rodapé
\usepackage{fancyhdr}

% Blocos
\usepackage[most]{tcolorbox}
\newtcolorbox{block}[1]{%
  colback=green!5,
  colframe=maincolor,
  coltitle=white,
  title=#1,
  fonttitle=\bfseries,
  boxrule=1.5pt,
  arc=3mm,
  left=6pt,
  right=6pt,
  top=6pt,
  bottom=6pt
}

% Margens
\geometry{margin=2.5cm}

% Fancyhdr
\pagestyle{fancy}
\fancyhf{}

\renewcommand{\headrulewidth}{2pt}
\renewcommand{\footrulewidth}{2pt}

\renewcommand{\headrule}{\hbox to\headwidth{%
  \color{forestgreen}\leaders\hrule height \headrulewidth\hfill}}

\renewcommand{\footrule}{\hbox to\headwidth{%
  \color{forestgreen}\leaders\hrule height \footrulewidth\hfill}}

% Altura do cabeçalho (corrigida)
\setlength{\headheight}{52pt}
\addtolength{\topmargin}{-20pt}

% Cabeçalho
\fancyhead[L]{%
  \includegraphics[height=1.2cm]{Universidade_de_Rio_Verde.jpg}
}

\fancyhead[R]{%
  \raggedleft
  \textcolor{forestgreen}{%
    \small
    \textbf{Lista de Exercícios}\\
    Me.~Sergio Souza Novak\\
    \textit{Banco de Dados}
    }
    }
    
% Rodapé
\fancyfoot[C]{\textcolor{forestgreen}{\thepage}}

\begin{document}

\vspace*{0.2cm}
\begin{center}
  {\Large \textbf{Lista de Exercícios: Normalização}}\\[0.2cm]
  {\large Banco de Dados}\\[0.2cm]
  Tema: 1ª e 2ª Formas Normais (1FN e 2FN)
\end{center}

\vspace{0.5cm}

\begin{enumerate}[label=\textbf{Questão \arabic*:}, leftmargin=*]

    \item \textbf{(Teoria - Anomalias)} Explique, com suas palavras, o que são anomalias de banco de dados e descreva brevemente os três tipos principais (Inserção, Atualização e Remoção).

    \item \textbf{(Anomalia de Inserção)} Considere a tabela abaixo, que armazena informações de turmas e disciplinas:
    
    \vspace{0.2cm}
    \centering
    \begin{tabular}{|l|l|l|l|}
        \hline
        \textbf{Cod\_Turma} & \textbf{Nome\_Turma} & \textbf{Cod\_Disc} & \textbf{Nome\_Disc} \\ \hline
        T1 & Turma A & D1 & Cálculo \\ \hline
        T2 & Turma B & D1 & Cálculo \\ \hline
    \end{tabular}
    \flushleft
    \vspace{0.2cm}
    Qual anomalia de inserção ocorre se tentarmos cadastrar uma nova disciplina (ex: "Física") que ainda não possui nenhuma turma vinculada? Como a normalização resolveria este problema?

    \item \textbf{(Anomalia de Atualização)} Na tabela \texttt{FUNCIONARIOS} abaixo, o departamento "Vendas" mudou seu nome para "Comercial".
    
    \vspace{0.2cm}
    \centering
    \begin{tabular}{|l|l|l|l|}
        \hline
        \textbf{Cod\_Fun} & \textbf{Nome} & \textbf{Cod\_Depto} & \textbf{Nome\_Depto} \\ \hline
        101 & Alice & D1 & Vendas \\ \hline
        102 & Bob & D1 & Vendas \\ \hline
        103 & Carol & D1 & Vendas \\ \hline
    \end{tabular}
    \flushleft
    \vspace{0.2cm}
    O que acontece se o DBA atualizar apenas o registro da Alice? Como chamamos esse problema e qual a causa raiz dele?

    \item \textbf{(1ª Forma Normal - Atributos Multivalorados)} A tabela abaixo viola a 1FN. Identifique o motivo e realize a normalização, separando os valores em novos registros.
    
    \vspace{0.2cm}
    \centering
    \begin{tabular}{|l|l|l|}
        \hline
        \textbf{ID} & \textbf{Nome} & \textbf{Contatos} \\ \hline
        1 & Carlos & car@email.com, (11) 9111-1111 \\ \hline
        2 & Ana & ana@email.com \\ \hline
        3 & Bia & bia@work.com, bia@home.com \\ \hline
    \end{tabular}
    \flushleft

    \item \textbf{(1ª Forma Normal - Grupos Repetitivos)} Um desenvolvedor criou a tabela \texttt{CLIENTES} com os campos: \texttt{ID, Nome, Tel\_1, Tel\_2, Tel\_3}.
    \begin{enumerate}[label=\alph*)]
        \item Por que essa estrutura viola a 1FN?
        \item Quais problemas essa modelagem pode causar no futuro?
        \item Refaça a modelagem utilizando uma nova tabela para armazenar os telefones.
    \end{enumerate}

    \item \textbf{(Dependência Funcional)} Dada a tabela \texttt{PRODUTOS}:
    
    \vspace{0.2cm}
    \centering
    \begin{tabular}{|l|l|l|l|}
        \hline
        \textbf{Cod\_Prod} & \textbf{Descricao} & \textbf{Cod\_Cat} & \textbf{Nome\_Categoria} \\ \hline
        P1 & Mouse & C1 & Periféricos \\ \hline
        P2 & Teclado & C1 & Periféricos \\ \hline
        P3 & Monitor & C2 & Vídeo \\ \hline
    \end{tabular}
    \flushleft
    \vspace{0.2cm}
    Identifique as Dependências Funcionais (DF) existentes. (Exemplo: $X \to Y$).

    \item \textbf{(2ª Forma Normal)} Para que uma tabela esteja na 2FN, ela deve obrigatoriamente satisfazer dois requisitos. Quais são eles? A 2FN é relevante para tabelas com chave primária simples (apenas um atributo)? Justifique.

    \item \textbf{(Normalização 2FN)} Considere a tabela de itens de um pedido, onde a Chave Primária é composta pelos campos \{\textbf{ID\_Pedido, ID\_Produto}\}:
    
    \vspace{0.2cm}
    \centering
    \begin{tabular}{|l|l|l|l|}
        \hline
        \textbf{ID\_Pedido} & \textbf{ID\_Produto} & \textbf{Quantidade} & \textbf{Nome\_Produto} \\ \hline
        501 & 10 & 2 & Notebook \\ \hline
        501 & 20 & 1 & Mouse \\ \hline
        502 & 10 & 5 & Notebook \\ \hline
    \end{tabular}
    \flushleft
    \vspace{0.2cm}
    \begin{enumerate}[label=\alph*)]
        \item Identifique a Dependência Funcional Parcial nesta tabela.
        \item Realize a decomposição da tabela para que ela passe a respeitar a 2FN.
    \end{enumerate}

\end{enumerate}

\vspace{1cm}
\begin{center}
    \textit{"A normalização é a arte de organizar os dados para que cada fato seja armazenado em apenas um lugar."}
\end{center}

\end{document}
